% =====================================================================
\lezione{02/10/2026}{Simulazione di laboratorio: severity assessment di un ransomware}
  {Prof.~Giuseppe Vaciago}
  {organizzazione dei laboratori e punti extra; caso Alpha S.r.l.\ (ransomware su piattaforma
   e-commerce B2B); ruoli e obblighi del responsabile; autovalutazione sul sito del Garante; modulo di
   notifica; metodologia ENISA di severity assessment ($SE = DPC \times EI + CB$); decisione sulla notifica}
\label{lez:2026-10-02}
% =====================================================================

\begin{nota}[Fonte]
Trascrizione automatica della registrazione, revisionata. Lezione \textbf{online}, svolta come
\textbf{simulazione} di laboratorio, con interazione via chat. Scale e soglie della metodologia sono
riprese dal documento \emph{03.~Data breach severity methodology} (ENISA, \emph{Recommendations for a
methodology of the assessment of severity of personal data breaches}, v1.0, dicembre 2013). In aula
si è usato anche lo strumento di autovalutazione sul sito del Garante (tradotto con il traduttore del
browser). I nomi delle aziende citate nel caso sono esempi di fantasia.
\end{nota}

\section{Sintesi}
Invece del laboratorio vero e proprio, il docente svolge una simulazione guidata su un caso
volutamente semplice: un ransomware blocca per 24 ore una piattaforma di e-commerce B2B che lavora per
conto di grandi marchi. Si ragiona su ruoli (la piattaforma è \textbf{responsabile}, i marchi sono
\textbf{titolari}), sul modulo di notifica al Garante e soprattutto sulla \textbf{metodologia ENISA di
severity assessment}, applicata passo per passo. Il risultato ($SE = 2$, rischio medio) porta alla
notifica al Garante; il docente consiglia di informare comunque gli utenti. Messaggio centrale: ogni
scelta nella formula pesa molto e va \textbf{motivata} (accountability); il CEO chiederà quasi sempre
di abbassare il punteggio.

% ---------------------------------------------------------------------
\section{Organizzazione}
\begin{itemize}
  \item \textbf{Calendario}: lunedì 5/10 lezione (in presenza) che chiude la parte di protezione dei
    dati con riferimento alla cybersicurezza; a fine lezione viene presentato il caso del primo
    \textbf{laboratorio ufficiale}, che si svolge \textbf{venerdì 9/10} online.
  \item \textbf{Punti extra}: \textbf{0,25 punti} sul voto d'esame per ogni laboratorio, per le
    elaborazioni migliori; non è garantito che tutti li ricevano. Le regole definitive saranno pubblicate
    dopo il primo laboratorio ufficiale, su una pagina dedicata.
  \item \textbf{Gruppi}: massimo 5 persone (6 tollerato, non 8--10), formati dagli studenti; il docente
    può aiutare chi non trova un gruppo. Si può lavorare anche da soli, ma il docente consiglia
    fortemente il gruppo.
  \item \textbf{Modalità}: il lab si svolge nell'ora e mezza del venerdì. I gruppi possono usare una
    propria call (Zoom, Meet) o trovarsi al Politecnico; per le domande il docente apre una stanza
    separata con il gruppo. Più tempo solo in casi eccezionali.
  \item Il documento ENISA viene caricato sulla piattaforma: \textbf{non è materia d'esame}, ma è ``la
    Bibbia'' per il laboratorio.
\end{itemize}

% ---------------------------------------------------------------------
\section{Il caso: Alpha S.r.l.}
\begin{itemize}
  \item \textbf{Alpha S.r.l.} gestisce in modalità SaaS l'e-commerce di circa \textbf{20 grandi
    clienti} (marchi di moda, bricolage\ldots): chi compra sul sito del marchio viene in realtà servito
    dalla piattaforma di Alpha, basata a sua volta su Shopify (un \textbf{sub-responsabile}). Alpha
    gestisce anche la logistica delle consegne e trattiene una percentuale sulle vendite.
  \item Un \textbf{ransomware} blocca tutte le transazioni per \textbf{24 ore}. Il team di risposta
    (DPO, legale, IT, CISO) verifica che i backup sono integri e, con i piani di \textbf{business
    continuity} e \textbf{disaster recovery}, ripristina tutto in 24 ore.
  \item Domanda: che cosa deve fare il DPO dal punto di vista della protezione dei dati? Per decidere se
    la notifica al Garante è obbligatoria serve un \textbf{severity assessment}.
\end{itemize}

\subsection{Primi passi}
\paragraph{Autovalutazione sul sito del Garante.} Lo strumento chiede: c'è stato un incidente di
sicurezza con perdita di riservatezza, integrità o disponibilità? Sì: almeno la
\textbf{disponibilità} è certa. Riguarda dati personali? Sì: i clienti di Alpha sono persone giuridiche,
ma ci sono i dati dei loro dipendenti e soprattutto quelli degli \textbf{utenti finali} che comprano
online. È quindi una violazione di dati personali. La pagina del Garante rimanda alla metodologia ENISA.

\paragraph{Chi notifica.} Alpha è \textbf{responsabile}: deve informare \textbf{senza ingiustificato
ritardo} ogni titolare, cioè ciascuno dei suoi clienti. Il paradosso tipico degli incidenti: notifica il
titolare, ma tutte le informazioni le ha il responsabile, che deve collaborare.
\begin{itemize}
  \item Con 20 clienti, \textbf{ogni titolare} fa la propria notifica (e l'autorità dovrà poi riunire
    20 pratiche sullo stesso caso).
  \item Eccezione pratica raccontata dal docente: in un breach di un'app di pagamento usata da moltissimi
    servizi, il Garante chiese al responsabile un'unica notifica per conto di tutti i titolari, per non
    riceverne un numero enorme. È una soluzione eccezionale, non la regola.
\end{itemize}

\begin{esame}
Ricordare ruoli e flusso: il \textbf{responsabile} informa il \textbf{titolare}; il \textbf{titolare}
notifica all'autorità. Nel caso, il titolare è il marchio, non la piattaforma.
\end{esame}

\subsection{Il modulo di notifica}
\begin{itemize}
  \item \textbf{Titolare}: il marchio cliente (si ipotizza che Alpha gli abbia girato tutte le
    informazioni).
  \item \textbf{Date}: data della violazione e data in cui il titolare ne ha avuto ragionevole certezza.
    È decisiva per le 72 ore. Nel caso: attacco il 30 settembre, consapevolezza il 1° ottobre. Spesso
    l'esfiltrazione precede di molto il ransomware ed è difficile datarla.
  \item \textbf{Natura}: disponibilità certa; la \textbf{riservatezza} è il vero nodo. Con la doppia
    estorsione l'attaccante chiede un riscatto e pubblica un campione dei dati esfiltrati sul dark web:
    a quel punto c'è una presunzione di perdita di riservatezza. A volte la causa si chiarisce solo dopo
    lunghe analisi.
  \item \textbf{Causa}: se non è nota, si dichiara che al momento non è nota.
  \item \textbf{Numero di interessati}: qualche centinaio di dipendenti dei clienti e, soprattutto, gli
    utenti: per la simulazione \textbf{50.000}.
  \item \textbf{Contatti del DPO}.
\end{itemize}

% ---------------------------------------------------------------------
\section{La metodologia ENISA}
\[
SE = DPC \times EI + CB
\]
\begin{itemize}
  \item \textbf{DPC} (\emph{Data Processing Context}): tipo di dati e contesto del trattamento.
  \item \textbf{EI} (\emph{Ease of Identification}): quanto è facile risalire alla persona.
  \item \textbf{CB} (\emph{Circumstances of the Breach}): perdita di riservatezza, integrità,
    disponibilità e intento malevolo; i punteggi si \textbf{sommano}.
\end{itemize}

\begin{center}
\small
\begin{tabularx}{\linewidth}{l X}
\toprule
\textbf{Fattore} & \textbf{Valori (punteggio di base)}\\ \midrule
DPC & dati semplici 1; comportamentali 2; finanziari 3; sensibili (artt.~9--10) 4\\
EI  & trascurabile 0,25; limitata 0,5; significativa 0,75; massima 1\\
CB -- riservatezza & 0 nessuna evidenza di trattamento illecito; +0,25 destinatari noti;
  +0,5 numero indeterminato di destinatari\\
CB -- integrità & 0; +0,25 alterazione recuperabile; +0,5 non recuperabile\\
CB -- disponibilità & 0 recupero senza difficoltà; +0,25 indisponibilità temporanea;
  +0,5 indisponibilità totale\\
CB -- intento malevolo & +0,5\\ \bottomrule
\end{tabularx}
\end{center}

\begin{center}
\small
\begin{tabular}{ll}
\toprule
\textbf{Punteggio} & \textbf{Severità}\\ \midrule
$SE < 2$ & bassa\\
$2 \le SE < 3$ & media\\
$3 \le SE < 4$ & alta\\
$SE \ge 4$ & molto alta\\ \bottomrule
\end{tabular}
\end{center}

\begin{esame}
La formula $SE = DPC \times EI + CB$ è ``la regola d'oro'': va \textbf{ricordata a memoria} e descritta
con chiarezza all'esame. È una ``formula stupida'' ma con conseguenze molto concrete.
\end{esame}

\subsection{Applicazione al caso}
\paragraph{DPC: quali dati?} Discussione con la classe:
\begin{itemize}
  \item \textbf{semplici} (nome, cognome, email, codice fiscale, telefono): base 1;
  \item \textbf{comportamentali}? Lo storico degli acquisti su più marchi permette di ricostruire un
    profilo (gusti, capacità di spesa). Il CEO direbbe di no, il DPO potrebbe dire di sì: è un punto su
    cui ``si può litigare'' a lungo;
  \item \textbf{finanziari}? Probabilmente no: i pagamenti passano spesso da servizi esterni (es.\
    Stripe) e il CEO afferma che i dati delle carte non sono conservati. Attenzione: per il GDPR i dati
    finanziari sono dati \emph{comuni}; è ENISA, con la sua logica di cybersicurezza, a pesarli 3;
  \item \textbf{sensibili}? Il data center potrebbe contenere dati sanitari dei dipendenti di Alpha, ma
    il titolare qui è il marchio, che non tratta dati di quel tipo.
\end{itemize}
Scelta del docente: \textbf{DPC = 1} (dati semplici), pur senza escludere i comportamentali, purché la
scelta sia motivata. ``Non ho una risposta: è la vostra scelta, il vostro livello di accountability.''

\paragraph{EI.} Nome, cognome, data di nascita, telefono insieme: \textbf{EI = 1}. Esempio opposto: un
solo indirizzo email composto da un numero di matricola renderebbe l'identificazione difficile
(0,25), mentre \texttt{nome.cognome@polito.it} no.

\paragraph{CB.}
\begin{itemize}
  \item \textbf{Riservatezza +0,25}: per ora i dati sono nelle mani dell'attaccante e forse sul dark
    web, ma non c'è certezza di una pubblicazione aperta. Per il docente il ``numero indeterminato di
    destinatari'' (+0,5) corrisponde a una pubblicazione accessibile a tutti; si potrebbe sostenere
    anche +0,5.
  \item \textbf{Integrità 0}: nessuna alterazione.
  \item \textbf{Disponibilità +0,25}: indisponibilità temporanea (24 ore).
  \item \textbf{Intento malevolo +0,5}: certamente sì.
\end{itemize}
\[
SE = 1 \times 1 + (0{,}25 + 0 + 0{,}25 + 0{,}5) = 2 \quad\Rightarrow\quad \text{media (``medio-bassa'')}
\]
Nella sezione sulle misure adottate conviene descrivere il ripristino in 24 ore: aiuta a giustificare
un punteggio contenuto.

\subsection{Dal punteggio alla decisione}
Lettura data in aula:
\begin{itemize}
  \item rischio \textbf{basso}: nessuna notifica al Garante;
  \item rischio \textbf{medio}: \textbf{notifica al Garante};
  \item rischio \textbf{alto o molto alto}: notifica al Garante e \textbf{comunicazione agli
    interessati}.
\end{itemize}
Con $SE = 2$ si è al confine: in teoria si potrebbe discutere se evitare la notifica, ma con tante
ipotesi incerte la strada sicura è notificare. Il docente consiglia anche di \textbf{avvisare gli
utenti}, pur non essendo obbligatorio, per creare un rapporto di fiducia e perché i dati potrebbero
essere usati per frodi. Se fossero stati coinvolti dati delle carte (DPC 3), il risultato sarebbe stato
alto o molto alto e la comunicazione agli interessati immediata.

\begin{integrazione}[le soglie del GDPR]
Nel GDPR la notifica al Garante (art.~33(1)) è dovuta salvo che sia \emph{improbabile} che la violazione
presenti un rischio per i diritti e le libertà delle persone; la comunicazione agli interessati
(art.~34(1)) è dovuta quando il rischio è \emph{elevato}. Le fasce ENISA sono uno strumento per
motivare queste due valutazioni, non la norma stessa.
\end{integrazione}

\subsection{Il dilemma del consulente}
Uno studente osserva che nessuno ha interesse a dare una valutazione peggiore del necessario, e che
quindi il punteggio verrà sempre abbassato. Il docente conferma: nel 95\% dei casi il CEO chiederà di
``ammorbidire'' la valutazione. Abbassare i punteggi è legittimo quando ha senso ed è motivato; farlo in
modo malevolo, nascondendo o falsando informazioni, espone alla sanzione (10 milioni o 2\% del
fatturato) e al rischio penale per false dichiarazioni al Garante (reclusione da sei mesi a tre anni),
che colpisce chi firma la dichiarazione. Sulla riservatezza: dati finiti sul dark web sono in generale
da considerare \textbf{persi}, anche se poi rimossi, ma la risposta va data caso per caso.

% ---------------------------------------------------------------------
\section{Domande aperte e punti da verificare}
\begin{verifica}
\begin{itemize}
  \item \textbf{Soglia di notifica}: nella parte finale la trascrizione è confusa (``obbligati con 2,5''?
    ``con 3''?). Qui si è usata la lettura esplicita data poco prima: basso $\rightarrow$ nessuna
    notifica; medio $\rightarrow$ Garante; alto/molto alto $\rightarrow$ anche interessati. Confrontare
    con la spiegazione teorica della lezione del 05/10.
  \item \textbf{Riservatezza +0,25 o +0,5}: nella metodologia ENISA +0,25 indica dati esposti a
    destinatari \emph{noti}; un attaccante che potrebbe rivendere i dati è difficilmente un destinatario
    noto. La scelta del docente è dichiaratamente discutibile.
  \item \textbf{Date del caso} (attacco il 30/09, consapevolezza il 1/10) e numero di utenti (50.000):
    ipotesi della simulazione, dette a voce.
  \item Il docente ha anche detto che la metodologia ENISA risale al 2014 (lezione del 28/09); il
    documento è datato dicembre 2013.
\end{itemize}
\end{verifica}
